% This document was generated by an LLM.

\documentclass{essay}

\usepackage{centernot}

\title{The First and Second Derivative Tests for Local Extrema}
\subtitle{A Real-Analysis Treatment}
\author{}
\date{}

\begin{document}

\maketitle
\tableofcontents

\section{Introduction}

Let \(D\subseteq\mathbb{R}\) and \(f:D\to\mathbb{R}\). This essay treats two
classical sufficient conditions for a differentiable function to have a
local extremum at a point --- the first and second derivative tests --- with
full statements of hypotheses and proofs, rather than as computational
recipes.

Stated as a recipe, the tests reduce to: solve \(f'(x)=0\), inspect \(f'\) or
\(f''\), and declare a maximum or minimum. This conceals several logical
subtleties that this essay makes explicit and proves:

\begin{enumerate}
    \item \(f'(c)=0\) is necessary for a differentiable interior extremum,
    but not sufficient.
    \item An extremum can occur where \(f'\) does not exist.
    \item The condition \(f''(c)>0\) implies a local minimum only under the
    additional hypothesis \(f'(c)=0\), and even then \(f''(c)=0\) settles
    nothing.
    \item Boundary points of the domain require separate treatment, since
    the tests are stated for interior points.
\end{enumerate}

We also develop concavity precisely, since it is the concept underlying the
informal language of a graph ``bending upward'' or ``bending downward'' used
to motivate the second derivative test.

\section{Local Extrema and Critical Points}

\begin{definition}[Local extremum]
Let \(f:D\to\mathbb{R}\) and \(c\in D\). We say \(f\) has a \emph{local
maximum} at \(c\) if
\[
\exists\,\delta>0 \quad \forall x\in D:\quad |x-c|<\delta \implies f(x)\leq f(c).
\]
We say \(f\) has a \emph{local minimum} at \(c\) if
\[
\exists\,\delta>0 \quad \forall x\in D:\quad |x-c|<\delta \implies f(x)\geq f(c).
\]
A \emph{local extremum} is a point at which \(f\) has a local maximum or a
local minimum. The extremum is \emph{strict} if the inequality is strict for
every \(x\neq c\) satisfying \(|x-c|<\delta\).
\end{definition}

\begin{definition}[Critical point]
A point \(c\in D\) is a \emph{critical point} of \(f\) if either \(f'(c)=0\)
or \(f'(c)\) fails to exist.
\end{definition}

Critical points are candidates for extrema; the definition asserts nothing
about whether \(c\) actually is one.

\begin{example}
Let \(f(x)=x^3\). Then \(f'(0)=0\), so \(0\) is a critical point, but for
every \(\delta>0\) there exist points of \((-\delta,0)\) with \(f<0=f(0)\)
and points of \((0,\delta)\) with \(f>0=f(0)\). Hence \(0\) is neither a
local maximum nor a local minimum.
\end{example}

\begin{example}
Let \(f(x)=|x|\). Then \(f'(0)\) does not exist, so \(0\) is a critical
point, and \(f(x)=|x|\geq 0=f(0)\) for every \(x\), so \(0\) is a (strict)
local minimum.
\end{example}

These two examples show that neither implication
\[
f'(c)=0 \implies c \text{ is an extremum}
\qquad\text{nor}\qquad
c \text{ is an extremum} \implies f'(c)=0
\]
holds in general; the first fails by \(x^3\), the second by \(|x|\).

\section{Fermat's Theorem: A Necessary Condition}

\begin{theorem}[Fermat]
Let \(f:D\to\mathbb{R}\), let \(c\) be an interior point of \(D\) (i.e.\
\((c-\varepsilon,c+\varepsilon)\subseteq D\) for some \(\varepsilon>0\)),
and suppose \(f\) is differentiable at \(c\) and has a local extremum at
\(c\). Then \(f'(c)=0\).
\end{theorem}

\begin{proof}
Suppose \(f\) has a local maximum at \(c\); the minimum case is symmetric.
Choose \(\delta\in(0,\varepsilon)\) such that \(f(x)\leq f(c)\) for all
\(x\) with \(|x-c|<\delta\). For \(x\in(c,c+\delta)\), both \(f(x)-f(c)\leq
0\) and \(x-c>0\), so
\[
\frac{f(x)-f(c)}{x-c}\leq 0.
\]
Since \(f\) is differentiable at \(c\), the one-sided limit as \(x\to
c^+\) exists and equals \(f'(c)\), so \(f'(c)\leq 0\).

For \(x\in(c-\delta,c)\), both \(f(x)-f(c)\leq0\) and \(x-c<0\), so
\[
\frac{f(x)-f(c)}{x-c}\geq 0,
\]
and letting \(x\to c^-\) gives \(f'(c)\geq 0\). Together, \(f'(c)\leq 0\)
and \(f'(c)\geq 0\) force \(f'(c)=0\).
\end{proof}

Both hypotheses are needed. Restricting \(f(x)=x\) to \(D=[0,1]\), the point
\(c=0\) is not interior to \(D\); \(f\) has a local minimum there, yet the
one-sided derivative at \(0\) is \(1\neq 0\). And \(f(x)=|x|\) has a local
minimum at the interior point \(0\), where \(f\) fails to be differentiable.

Fermat's theorem gives a necessary condition,
\[
\text{differentiable interior extremum} \implies f'(c)=0,
\]
and the converse is false:
\[
f'(c)=0 \quad\centernot\Longrightarrow\quad c \text{ is a local extremum}.
\]
The derivative tests below supply additional hypotheses under which a
critical point \emph{can} be classified.

\section{The First Derivative Test}

\begin{theorem}[First derivative test]
Let \(f:D\to\mathbb{R}\), let \(c\) be an interior point of \(D\), and
suppose \(f\) is continuous at \(c\) and there exists \(\delta>0\) such that
\(f\) is differentiable on \((c-\delta,c)\cup(c,c+\delta)\).
\begin{enumerate}
    \item If \(f'<0\) on \((c-\delta,c)\) and \(f'>0\) on \((c,c+\delta)\),
    then \(f\) has a strict local minimum at \(c\).
    \item If \(f'>0\) on \((c-\delta,c)\) and \(f'<0\) on \((c,c+\delta)\),
    then \(f\) has a strict local maximum at \(c\).
    \item If \(f'\) has the same sign throughout \((c-\delta,c)\cup(c,c+\delta)\),
    then \(f\) has no local extremum at \(c\).
\end{enumerate}
\end{theorem}

\begin{proof}
We prove (1); the proof of (2) is obtained by replacing \(f\) with \(-f\).

Let \(x\in(c-\delta,c)\). Since \(f\) is differentiable on \((c-\delta,c)\)
and continuous at \(c\), \(f\) is continuous on \([x,c]\) and differentiable
on \((x,c)\), so the Mean Value Theorem gives \(\xi\in(x,c)\) with
\[
\frac{f(c)-f(x)}{c-x}=f'(\xi).
\]
Since \(\xi\in(c-\delta,c)\), \(f'(\xi)<0\), and \(c-x>0\), so
\(f(c)-f(x)<0\), i.e.\ \(f(x)>f(c)\).

Let \(x\in(c,c+\delta)\). Likewise the Mean Value Theorem on \([c,x]\) gives
\(\eta\in(c,x)\) with
\[
\frac{f(x)-f(c)}{x-c}=f'(\eta)>0,
\]
so \(f(x)>f(c)\).

Hence \(f(x)>f(c)\) for every \(x\in(c-\delta,c+\delta)\setminus\{c\}\), which
is a strict local minimum.

For (3), suppose \(f'>0\) throughout \((c-\delta,c)\cup(c,c+\delta)\) (the
case \(f'<0\) is symmetric). The same Mean Value Theorem argument on
\([x,c]\) or \([c,x]\), for \(x\) on either side of \(c\), shows \(f\) is
strictly increasing on \((c-\delta,c+\delta)\); in particular \(f(x)<f(c)\)
for \(x<c\) and \(f(x)>f(c)\) for \(x>c\) in this interval, so \(c\) is
neither a local maximum nor a local minimum.
\end{proof}

The hypothesis doing the real work is not \(f'(c)=0\) but a \emph{sign
change} of \(f'\) across \(c\); the theorem does not even require \(f\) to
be differentiable at \(c\) itself, only continuous there.

\section{Examples of the First Derivative Test}

\begin{example}[Local minimum]
For \(f(x)=x^2\), \(f'(x)=2x\), which is negative on \((-\infty,0)\) and
positive on \((0,\infty)\). By the first derivative test, \(0\) is a strict
local minimum.
\end{example}

\begin{example}[Local maximum]
For \(f(x)=-x^2\), \(f'(x)=-2x\) is positive on \((-\infty,0)\) and negative
on \((0,\infty)\), so \(0\) is a strict local maximum.
\end{example}

\begin{example}[Critical point with no sign change]
For \(f(x)=x^3\), \(f'(x)=3x^2>0\) for every \(x\neq 0\); the sign does not
change across \(0\), so \(0\) is not a local extremum, consistent with
\S 2.
\end{example}

\begin{example}[Extremum at a point of non-differentiability]
For \(f(x)=|x|\), \(f'(x)=-1\) on \((-\infty,0)\) and \(f'(x)=1\) on
\((0,\infty)\); \(f\) is continuous at \(0\) although \(f'(0)\) does not
exist. The hypotheses of the first derivative test are satisfied with
\(c=0\), and the sign change gives a local minimum. This is the key
advantage of the first derivative test over the second: it does not require
differentiability at \(c\).
\end{example}

\section{Limits of the First Derivative Test}

\begin{remark}[No fixed-sign neighborhood]
The theorem requires the existence of a single \(\delta>0\) on which \(f'\)
has constant sign on each side of \(c\). This need not hold. Let
\[
f(x)=\begin{cases} x^4\sin(1/x), & x\neq 0,\\ 0, & x=0.\end{cases}
\]
One computes \(f'(x)=4x^3\sin(1/x)-x^2\cos(1/x)\) for \(x\neq0\), which
changes sign infinitely often on every interval \((0,\delta)\) as
\(1/x\to\infty\). No \(\delta>0\) makes \(f'\) of constant sign on
\((0,\delta)\), so the theorem's hypotheses are not met at \(c=0\) and the
sign-chart method gives no direct conclusion, even though \(0\) is in fact
a strict local minimum here (since \(|f(x)|\le x^4\)).
\end{remark}

\begin{remark}[Continuity at \(c\) cannot be dropped]
The hypothesis that \(f\) be continuous \emph{at} \(c\) --- not merely
differentiable on the punctured neighborhood --- is essential, since the
theorem's conclusion is about the value \(f(c)\) itself. Let
\[
f(x)=\begin{cases} x^2, & x\neq 0,\\ 1, & x=0.\end{cases}
\]
Here \(f'(x)=2x\) is negative for \(x<0\) and positive for \(x>0\), the same
pattern as in the local-minimum case of the theorem. But \(f\) is not
continuous at \(0\): \(f(0)=1\) while \(f(x)\to 0\) as \(x\to 0\).
Consequently \(f(x)<f(0)\) for all \(x\neq0\) near \(0\), so \(0\) is
actually a strict local \emph{maximum}. The sign pattern of \(f'\) on the
punctured neighborhood constrains only the relative order of \(f\) on the
two sides; continuity at \(c\) is what ties that behavior to \(f(c)\).
\end{remark}

\begin{remark}[Endpoints]
The theorem is stated for interior points and gives no information at
endpoints, where \(f'\) is defined (if at all) on only one side. For
\(f(x)=x\) on \([0,1]\), \(f\) has a minimum at \(0\) and a maximum at
\(1\), with no sign change of \(f'\equiv1\) available at either point.
Endpoints must be examined by direct comparison of values, not by the sign
of \(f'\) on two sides.
\end{remark}

\section{Concavity}

\begin{definition}[Convex and concave functions]
Let \(I\subseteq\mathbb{R}\) be an interval and \(f:I\to\mathbb{R}\). We say
\(f\) is \emph{convex} (or \emph{concave up}) on \(I\) if
\[
f(tx+(1-t)y)\leq t f(x)+(1-t)f(y)
\qquad \forall x,y\in I,\ \forall t\in[0,1].
\]
We say \(f\) is \emph{concave} (or \emph{concave down}) on \(I\) if the
reverse inequality holds. The inequalities are non-strict, so a linear
function is both convex and concave.
\end{definition}

\begin{proposition}[Monotone derivative criterion]
Let \(f\) be differentiable on an interval \(I\). Then \(f\) is convex on
\(I\) if and only if \(f'\) is nondecreasing on \(I\).
\end{proposition}

\begin{proof}
(\(\Leftarrow\)) Suppose \(f'\) is nondecreasing. Fix \(x<y\) in \(I\) and
\(t\in(0,1)\), and let \(z=tx+(1-t)y\), so \(x<z<y\). By the Mean Value
Theorem there exist \(\xi_1\in(x,z)\) and \(\xi_2\in(z,y)\) with
\[
\frac{f(z)-f(x)}{z-x}=f'(\xi_1), \qquad \frac{f(y)-f(z)}{y-z}=f'(\xi_2).
\]
Since \(\xi_1<z<\xi_2\) and \(f'\) is nondecreasing, \(f'(\xi_1)\leq
f'(\xi_2)\), i.e.
\[
\frac{f(z)-f(x)}{z-x}\leq\frac{f(y)-f(z)}{y-z}.
\]
Using \(z-x=(1-t)(y-x)\) and \(y-z=t(y-x)\) and clearing denominators
(both positive) yields, after rearranging,
\[
f(z)\leq t f(x)+(1-t)f(y),
\]
which is the convexity inequality at \(z=tx+(1-t)y\).

(\(\Rightarrow\)) Suppose \(f\) is convex and let \(x<y\) in \(I\). A
standard consequence of the convexity inequality is that the difference
quotient \(\bigl(f(w)-f(x)\bigr)/(w-x)\) is nondecreasing in \(w>x\); letting
\(w\to x^+\) and \(w\to y^-\) shows \(f'(x)\leq f'(y)\).
\end{proof}

\begin{proposition}[Second-derivative criterion]
If \(f\) is twice differentiable on an interval \(I\) and \(f''\geq0\) on
\(I\), then \(f\) is convex on \(I\); if \(f''\leq0\) on \(I\), then \(f\)
is concave on \(I\).
\end{proposition}

\begin{proof}
If \(f''\geq0\) on \(I\), then for \(x<y\) in \(I\) the Mean Value Theorem
applied to \(f'\) gives \(\xi\in(x,y)\) with \(f'(y)-f'(x)=f''(\xi)(y-x)\geq
0\), so \(f'\) is nondecreasing on \(I\). Convexity follows from the
previous proposition. The concave case is analogous.
\end{proof}

\begin{proposition}[Tangent-line characterization]
If \(f\) is differentiable and convex on \(I\), then for all \(a,x\in I\),
\[
f(x)\geq f(a)+f'(a)(x-a).
\]
If \(f\) is differentiable and concave on \(I\), the reverse inequality
holds.
\end{proposition}

\begin{proof}
Suppose \(f\) is convex, so \(f'\) is nondecreasing by the monotone
derivative criterion. For \(x>a\), the Mean Value Theorem gives \(\xi\in
(a,x)\) with \((f(x)-f(a))/(x-a)=f'(\xi)\geq f'(a)\), so \(f(x)-f(a)\geq
f'(a)(x-a)\). For \(x<a\), the Mean Value Theorem on \([x,a]\) gives
\(\xi\in(x,a)\) with \((f(a)-f(x))/(a-x)=f'(\xi)\leq f'(a)\); multiplying by
\(a-x>0\) and rearranging again gives \(f(x)\geq f(a)+f'(a)(x-a)\). The case
\(x=a\) is trivial. The concave case follows by applying this to \(-f\).
\end{proof}

The geometric content is that the graph of a convex function lies on or
above every tangent line, and the graph of a concave function lies on or
below every tangent line.

\begin{example}
\(f(x)=x^2\) has \(f''(x)=2>0\), so \(f\) is convex on \(\mathbb{R}\); its
graph lies above every tangent line. \(f(x)=-x^2\) has \(f''(x)=-2<0\), so
\(f\) is concave on \(\mathbb{R}\).
\end{example}

\begin{definition}[Inflection point]
A point \(c\) is an \emph{inflection point} of \(f\) if \(f\) is convex on
one side of \(c\) and concave on the other.
\end{definition}

\begin{example}
For \(f(x)=x^3\), \(f''(x)=6x\), which is negative for \(x<0\) and positive
for \(x>0\). So \(f\) is concave on \((-\infty,0)\) and convex on
\((0,\infty)\), and \(0\) is an inflection point. Note that \(0\) is not a
local extremum of \(x^3\): a change of concavity is a distinct phenomenon
from an extremum.
\end{example}

\section{The Second Derivative Test}

\begin{theorem}[Second derivative test]
Let \(c\) be an interior point of the domain of \(f\), suppose \(f'\) is
defined on a neighborhood of \(c\) with \(f'(c)=0\), and suppose \(f''(c)\)
exists.
\begin{enumerate}
    \item If \(f''(c)>0\), then \(f\) has a strict local minimum at \(c\).
    \item If \(f''(c)<0\), then \(f\) has a strict local maximum at \(c\).
    \item If \(f''(c)=0\), the test is inconclusive.
\end{enumerate}
\end{theorem}

\begin{proof}
Suppose \(f''(c)>0\). By definition, and using \(f'(c)=0\),
\[
f''(c)=\lim_{x\to c}\frac{f'(x)-f'(c)}{x-c}=\lim_{x\to c}\frac{f'(x)}{x-c}.
\]
Since this limit is positive, there exists \(\delta>0\) such that
\(f'(x)/(x-c)>0\) for all \(x\) with \(0<|x-c|<\delta\). For \(x\in
(c-\delta,c)\), \(x-c<0\), forcing \(f'(x)<0\); for \(x\in(c,c+\delta)\),
\(x-c>0\), forcing \(f'(x)>0\). This is exactly the sign pattern
\(f'<0\) on \((c-\delta,c)\) and \(f'>0\) on \((c,c+\delta)\), so the first
derivative test gives a strict local minimum at \(c\). The case
\(f''(c)<0\) is symmetric.

For (3), \(x^4\), \(-x^4\), and \(x^3\) (\S\S 12) all satisfy \(f'(0)=0=f''(0)\)
while exhibiting a strict local minimum, a strict local maximum, and no
extremum respectively; hence no conclusion of the stated form can hold when
\(f''(c)=0\).
\end{proof}

The proof shows the second derivative test is not an independent
phenomenon: a nonzero \(f''(c)\), together with \(f'(c)=0\), forces exactly
the sign change of \(f'\) required by the first derivative test.

\begin{remark}
The stationary-point hypothesis \(f'(c)=0\) cannot be dropped. For
\(f(x)=x^2+x\), \(f''(0)=2>0\), but \(f'(0)=1\neq0\) and \(0\) is not a
local extremum (the minimum of this function is at \(x=-1/2\)). So
\[
f''(c)>0 \quad\centernot\Longrightarrow\quad c \text{ is a local minimum}
\]
without the hypothesis \(f'(c)=0\).
\end{remark}

\section{Taylor's Theorem as a Local Model}

\begin{definition}[Little-\(o\)]
For a function \(g\) defined near \(0\) and \(n\in\mathbb{N}\), we write
\(g(h)=o(h^n)\) as \(h\to0\) if
\[
\lim_{h\to0}\frac{g(h)}{h^n}=0.
\]
\end{definition}

\begin{theorem}[Taylor, Peano form]
Suppose \(f^{(n-1)}\) exists on a neighborhood of \(c\) and \(f^{(n)}(c)\)
exists. Then
\[
f(c+h)=f(c)+f'(c)h+\frac{f''(c)}{2!}h^2+\cdots+\frac{f^{(n)}(c)}{n!}h^n+o(h^n)
\qquad\text{as } h\to0.
\]
\end{theorem}

We take this theorem as known; it is a standard result of one-variable
calculus, proved by repeated application of L'H\^opital's rule or Cauchy's
mean value theorem to the remainder. We use it to re-derive the two
derivative tests as corollaries, with the sign arguments made fully
rigorous via the definition of \(o(h^n)\).

\begin{corollary}[Fermat's theorem, again]
If \(f'(c)=0\) fails, i.e.\ \(f'(c)\neq0\), then \(c\) is not an interior
local extremum.
\end{corollary}

\begin{proof}
By Taylor's theorem with \(n=1\), \(f(c+h)-f(c)=f'(c)h+o(h)\). Let
\(\varepsilon=|f'(c)|/2>0\); there is \(\delta>0\) such that
\(|o(h)|<\varepsilon|h|\) for \(0<|h|<\delta\). Then for such \(h\),
\(f(c+h)-f(c)\) has the same sign as \(f'(c)h\), namely the sign of
\(f'(c)\) when \(h>0\) and the opposite sign when \(h<0\). Since these
signs differ, \(f(c+h)-f(c)\) takes both signs on every punctured
neighborhood of \(c\), so \(c\) is not a local extremum.
\end{proof}

\begin{corollary}[Second derivative test, again]
Suppose \(f'(c)=0\) and \(f''(c)\) exists. If \(f''(c)>0\), \(c\) is a
strict local minimum; if \(f''(c)<0\), \(c\) is a strict local maximum.
\end{corollary}

\begin{proof}
By Taylor's theorem with \(n=2\) and \(f'(c)=0\),
\[
f(c+h)-f(c)=\frac{f''(c)}{2}h^2+o(h^2).
\]
Suppose \(f''(c)>0\) and let \(\varepsilon=f''(c)/4>0\). Choose \(\delta>0\)
such that \(|o(h^2)|<\varepsilon h^2\) for \(0<|h|<\delta\). Then for such
\(h\),
\[
f(c+h)-f(c) \geq \frac{f''(c)}{2}h^2-\varepsilon h^2 = \frac{f''(c)}{4}h^2>0.
\]
So \(f(c+h)>f(c)\) throughout the punctured neighborhood \(0<|h|<\delta\),
a strict local minimum. The case \(f''(c)<0\) is symmetric.
\end{proof}

Unlike the earlier proof from the limit definition of \(f''(c)\), this
argument makes the role of the remainder explicit: the conclusion depends
on the \(o(h^2)\) term being eventually dominated by the quadratic term,
which is exactly the content of the little-\(o\) definition.

\section{Examples of the Second Derivative Test}

\begin{example}
For \(f(x)=x^2+3\), \(f'(x)=2x\) and \(f''(x)=2\). At \(c=0\), \(f'(0)=0\)
and \(f''(0)=2>0\), so \(0\) is a strict local minimum.
\end{example}

\begin{example}
For \(f(x)=5-x^2\), \(f'(x)=-2x\) and \(f''(x)=-2\). At \(c=0\), \(f'(0)=0\)
and \(f''(0)=-2<0\), so \(0\) is a strict local maximum.
\end{example}

\section{The Case \texorpdfstring{\(f''(c)=0\)}{f''(c)=0}}

When \(f'(c)=0=f''(c)\), all three outcomes are possible.

\begin{example}
\(f(x)=x^4\) has \(f'(0)=f''(0)=0\), and \(x^4\geq0=f(0)\) for all \(x\), so
\(0\) is a strict local minimum.
\end{example}

\begin{example}
\(f(x)=-x^4\) has \(f'(0)=f''(0)=0\) and \(0\) is a strict local maximum, by
the same argument applied to \(-x^4\).
\end{example}

\begin{example}
\(f(x)=x^3\) has \(f'(0)=f''(0)=0\), yet \(0\) is neither a maximum nor a
minimum (\S 2); it is instead an inflection point, since \(f''(x)=6x\)
changes sign at \(0\) (\S 7).
\end{example}

Since all three outcomes occur under identical data \(f'(c)=f''(c)=0\), no
theorem of the form ``\(f'(c)=0=f''(c)\) implies \dots'' concerning the
extremum classification can be true; this is precisely the content of
``inconclusive.''

\section{An Extremum with No Useful Higher-Order Data}

\begin{example}
\(f(x)=|x|\) has a strict local minimum at \(0\), but neither \(f'(0)\) nor
\(f''(0)\) exists: the second derivative test is a sufficient condition,
never a necessary one.
\end{example}

\begin{example}
\(f(x)=|x|^{3/2}\) is differentiable at \(0\) with \(f'(0)=0\)
(indeed \(f'(x)=\tfrac32\operatorname{sgn}(x)\sqrt{|x|}\to0\) as
\(x\to0\)), and \(0\) is a strict local minimum since
\(|x|^{3/2}\geq0=f(0)\). But
\(f''(x)=\tfrac34|x|^{-1/2}\to\infty\) as \(x\to0\), so \(f''(0)\) does not
exist and the second derivative test cannot be applied.
\end{example}

\section{Necessary versus Sufficient Conditions}

For an interior point \(c\) at which \(f\) is differentiable, Fermat's
theorem gives the necessary condition
\[
c \text{ is a local extremum} \implies f'(c)=0,
\]
which is not sufficient (\(x^3\)). Under the stated neighborhood
hypotheses, a sign change of \(f'\) across \(c\) is sufficient for an
extremum (first derivative test), and \(f'(c)=0\) together with
\(f''(c)\neq0\) is sufficient (second derivative test), but neither test
is necessary: an extremum can occur where \(f'\) fails to exist (\(|x|\))
or where \(f''\) fails to exist or vanishes (\(x^4\)). The logical
structure is asymmetric: differentiable extrema are always critical points,
but not every classification tool applies to every critical point.

\section{Strict versus Non-Strict Extrema}

The derivative tests as stated detect \emph{strict} extrema, via an actual
sign change of \(f'\). Non-strict extrema need not exhibit such a change.

\begin{example}
The constant function \(f\equiv0\) has \(f(x)=f(c)\) for all \(x,c\), so
every point is simultaneously a (non-strict) local maximum and minimum, yet
\(f'\equiv0\) exhibits no sign change anywhere.
\end{example}

\begin{example}
Let
\[
f(x)=\begin{cases} x^2, & x<0,\\ 0, & 0\leq x\leq1,\\ (x-1)^2,& x>1.\end{cases}
\]
Every point of \([0,1]\) is a (non-strict) local minimum, but an interior
point of \((0,1)\) has \(f'\equiv0\) on a full neighborhood, not a sign
change from negative to positive.
\end{example}

The first derivative test, as stated, is a sufficient condition for a
\emph{strict} extremum; it should not be read as characterizing every
non-strict extremum.

\section{Endpoints and Constrained Domains}

For \(f:[a,b]\to\mathbb{R}\), an absolute extremum can occur at an interior
critical point, at an interior point of non-differentiability, or at either
endpoint \(a\) or \(b\) --- points to which Fermat's theorem and the two
derivative tests, stated for interior points, do not apply. The standard
method for locating an absolute extremum on \([a,b]\) is therefore to
evaluate \(f\) at every interior critical point and at both endpoints, and
compare the resulting values directly, rather than to rely on derivative
sign patterns at the endpoints.

\begin{example}
\(f(x)=x\) on \([0,1]\) has no interior critical point, yet
\(\min_{[0,1]}f=f(0)=0\) and \(\max_{[0,1]}f=f(1)=1\), found by direct
comparison of the two endpoint values.
\end{example}

\section{Higher-Order Derivative Tests}

\begin{theorem}[\(n\)th derivative test]
Suppose \(f'(c)=f''(c)=\cdots=f^{(n-1)}(c)=0\) and \(f^{(n)}(c)\neq0\),
where \(f^{(n-1)}\) exists on a neighborhood of \(c\).
\begin{enumerate}
    \item If \(n\) is even and \(f^{(n)}(c)>0\), \(c\) is a strict local
    minimum.
    \item If \(n\) is even and \(f^{(n)}(c)<0\), \(c\) is a strict local
    maximum.
    \item If \(n\) is odd, \(c\) is not a local extremum.
\end{enumerate}
\end{theorem}

\begin{proof}
By Taylor's theorem (Peano form), since all derivatives up to order
\(n-1\) vanish at \(c\),
\[
f(c+h)-f(c)=\frac{f^{(n)}(c)}{n!}h^n+o(h^n).
\]
Let \(\varepsilon=|f^{(n)}(c)|/(2\cdot n!)\); choose \(\delta>0\) such that
\(|o(h^n)|<\varepsilon|h|^n\) for \(0<|h|<\delta\). For such \(h\),
\(f(c+h)-f(c)\) has the same sign as \(f^{(n)}(c)h^n/n!\).

If \(n\) is even, \(h^n>0\) for every \(h\neq0\), so this sign equals the
sign of \(f^{(n)}(c)\) on both sides of \(c\), giving (1) and (2) exactly
as in the second-derivative case.

If \(n\) is odd, \(h^n\) has the sign of \(h\), so \(f(c+h)-f(c)\) has
opposite signs for \(h>0\) and \(h<0\); hence \(c\) is not a local
extremum.
\end{proof}

\begin{example}
For \(f(x)=x^4\), the first nonvanishing derivative at \(0\) is
\(f^{(4)}(0)=24>0\) with \(n=4\) even, giving a local minimum, matching
\S 12.
\end{example}

\begin{example}
For \(f(x)=x^3\), the first nonvanishing derivative at \(0\) is
\(f'''(0)=6\) with \(n=3\) odd, so \(0\) is not a local extremum, matching
\S 2.
\end{example}

\section{Smooth Is Not Analytic: A Counterexample}

The \(n\)th derivative test presupposes some \(n\) with \(f^{(n)}(c)\neq0\).
This can fail even for an infinitely differentiable function.

\begin{example}
Let
\[
f(x)=\begin{cases}e^{-1/x^2}, & x\neq0,\\0,& x=0.\end{cases}
\]
A standard computation shows \(f\) is infinitely differentiable at \(0\)
with \(f^{(n)}(0)=0\) for every \(n\geq1\). Yet \(f(x)>0=f(0)\) for every
\(x\neq0\), so \(0\) is a strict local minimum. No finite-order derivative
test can detect this, since every derivative at \(0\) vanishes: the Taylor
series of \(f\) at \(0\) is identically \(0\), while \(f\) itself is not
identically \(0\) on any neighborhood of \(0\). This is the standard
example separating smooth (\(C^\infty\)) functions from analytic ones.
\end{example}

\section{Comparison of the Two Tests}

The first derivative test uses information about \(f'\) on a punctured
neighborhood of \(c\) and requires only continuity, not differentiability,
at \(c\) itself; the second derivative test uses only the two numbers
\(f'(c)\) and \(f''(c)\), which is a strictly stronger requirement but
often faster to check. Consequently there are points the first derivative
test classifies that the second cannot even address: \(f(x)=x^4\) is
classified by the second derivative test only after it fails (\(f''(0)=0\))
and one falls back to a higher-order or first-derivative argument, while
\(f(x)=|x|\) is classified by the first derivative test even though
\(f''(0)\) is not defined.

\section{A Logical Checklist for Extremum Problems}

\begin{enumerate}
    \item Determine the domain, including endpoints and excluded points,
    before differentiating.
    \item Locate the candidate points: solutions of \(f'(x)=0\), points
    where \(f'\) fails to exist, and (for absolute extrema) the endpoints.
    \item At each interior candidate, check whether the first derivative
    test's hypotheses hold, i.e.\ whether \(f'\) has a fixed sign on each
    side of \(c\).
    \item If \(f'(c)=0\), compute \(f''(c)\); a nonzero value classifies
    \(c\) immediately by the second derivative test.
    \item If \(f''(c)=0\), draw no conclusion from this alone; return to
    the sign of \(f'\) directly, or to a higher-order derivative test if
    its hypotheses hold.
    \item For an absolute extremum on a closed interval, compare the
    values of \(f\) at all candidates and both endpoints directly.
\end{enumerate}

\section{Common Logical Errors}

\begin{enumerate}
    \item \emph{``\(f'(c)=0\) implies \(c\) is an extremum.''} False;
    \(f(x)=x^3\) at \(c=0\).
    \item \emph{``\(f''(c)>0\) implies \(c\) is a minimum.''} False without
    \(f'(c)=0\); \(f(x)=x^2+x\) at \(c=0\).
    \item \emph{``\(f''(c)=0\) implies \(c\) is an inflection point.''}
    False; \(f(x)=x^4\) at \(c=0\) has \(f''(0)=0\) but no change of
    concavity there, since \(f''(x)=12x^2\geq0\) on both sides.
    \item \emph{``If the second derivative test fails, there is no
    extremum.''} False; \(f(x)=x^4\) at \(c=0\).
    \item \emph{``Every extremum occurs where \(f'(c)=0\).''} False;
    \(f(x)=|x|\) at \(c=0\), and any endpoint extremum.
    \item \emph{``A critical point is an extremum.''} False in general; a
    critical point is only a candidate.
\end{enumerate}

\section{Summary}

\[
\begin{array}{c|c|c}
\text{Hypothesis} & \text{Conclusion} & \text{Source}\\
\hline
f'(c)=0 & \text{candidate only} & \text{Fermat (necessary, not sufficient)}\\[4pt]
f':-\to+ \text{ across } c & \text{strict local minimum} & \text{first derivative test}\\[4pt]
f':+\to- \text{ across } c & \text{strict local maximum} & \text{first derivative test}\\[4pt]
f' \text{ fixed sign across } c & \text{no local extremum} & \text{first derivative test}\\[4pt]
f''\geq0 \text{ on } I & f \text{ convex on } I & \text{concavity criterion}\\[4pt]
f''\leq0 \text{ on } I & f \text{ concave on } I & \text{concavity criterion}\\[4pt]
f'' \text{ changes sign across } c & \text{inflection point} & \text{concavity change}\\[4pt]
f'(c)=0,\ f''(c)>0 & \text{strict local minimum} & \text{second derivative test}\\[4pt]
f'(c)=0,\ f''(c)<0 & \text{strict local maximum} & \text{second derivative test}\\[4pt]
f'(c)=0,\ f''(c)=0 & \text{inconclusive} & ---
\end{array}
\]

\section{Conclusion}

The derivative tests are theorems, not definitions: they relate local sign
and vanishing behavior of derivatives to the local order structure of
\(f\), under hypotheses that must be verified rather than assumed. The
first derivative test is fundamentally a statement about monotonicity,
proved from the Mean Value Theorem; the second derivative test is a
sufficient condition, at a stationary point, that forces the sign change
required by the first, provable either from the limit definition of
\(f''(c)\) or from Taylor's theorem with Peano remainder. Concavity is a
logically separate notion --- monotonicity of \(f'\) rather than sign of
\(f\) or \(f'\) --- and an inflection point is a change of concavity, not
an extremum.

Every implication proved here runs in one direction only. \(f'(c)=0,\
f''(c)>0\) is sufficient, not necessary, for a strict local minimum: a
minimum may occur with \(f''(c)=0\) (\(x^4\)), with \(f''(c)\) undefined
(\(|x|^{3/2}\)), with \(f'(c)\) undefined (\(|x|\)), or with every
derivative at \(c\) vanishing while \(f\) is still smooth
(\(e^{-1/x^2}\)). Correct use of these tests therefore requires checking,
at each step, exactly which hypothesis is being invoked and whether it
actually holds --- not merely computing derivatives and reading off a sign.

\end{document}
